% >>> typage-template pkuthss@2026.10.1 variant=standard
\documentclass[UTF8,fontset=fandol]{pkuthss}
\usepackage[backend=biber,style=caspervector,utf8,sorting=ecnyt]{biblatex}
% Double-blind version: the "blind" variant adds \blindtrue below.
\newif\ifblind
% <<< typage-template
\pkuthssinfo{
  cthesisname = {博士研究生学位论文}, ethesisname = {Doctoral Dissertation},
  thesiscover = {博士研究生学位论文},
  ctitle = {论文题目}, etitle = {Thesis Title},
  cauthor = {作者姓名}, eauthor = {Author Name},
  date = {二〇二七年六月},
  studentid = {0000000000}, school = {某某学院},
  cmajor = {某某专业}, emajor = {Major},
  direction = {某某方向},
  cmentor = {导师姓名\ 教授}, ementor = {Prof.\ Supervisor Name},
  ckeywords = {关键词一，关键词二，关键词三},
  ekeywords = {keyword one, keyword two, keyword three},
  degreetype = {1},
  blindid = {0000000000}, discipline = {某某学科},
}
\addbibresource{refs.bib}
\usepackage{booktabs}

\begin{document}

\frontmatter
\pagestyle{empty}
\ifblind\makeblind\else\maketitle\fi

\cleardoublepage
% Copyright (c) 2008-2009 solvethis
% Copyright (c) 2010-2017,2021 Casper Ti. Vector
% All rights reserved.
%
% Redistribution and use in source and binary forms, with or without
% modification, are permitted provided that the following conditions are
% met:
%
% * Redistributions of source code must retain the above copyright notice,
%   this list of conditions and the following disclaimer.
% * Redistributions in binary form must reproduce the above copyright
%   notice, this list of conditions and the following disclaimer in the
%   documentation and/or other materials provided with the distribution.
% * Neither the name of Peking University nor the names of its contributors
%   may be used to endorse or promote products derived from this software
%   without specific prior written permission.
%
% THIS SOFTWARE IS PROVIDED BY THE COPYRIGHT HOLDERS AND CONTRIBUTORS "AS
% IS" AND ANY EXPRESS OR IMPLIED WARRANTIES, INCLUDING, BUT NOT LIMITED TO,
% THE IMPLIED WARRANTIES OF MERCHANTABILITY AND FITNESS FOR A PARTICULAR
% PURPOSE ARE DISCLAIMED. IN NO EVENT SHALL THE COPYRIGHT HOLDER OR
% CONTRIBUTORS BE LIABLE FOR ANY DIRECT, INDIRECT, INCIDENTAL, SPECIAL,
% EXEMPLARY, OR CONSEQUENTIAL DAMAGES (INCLUDING, BUT NOT LIMITED TO,
% PROCUREMENT OF SUBSTITUTE GOODS OR SERVICES; LOSS OF USE, DATA, OR
% PROFITS; OR BUSINESS INTERRUPTION) HOWEVER CAUSED AND ON ANY THEORY OF
% LIABILITY, WHETHER IN CONTRACT, STRICT LIABILITY, OR TORT (INCLUDING
% NEGLIGENCE OR OTHERWISE) ARISING IN ANY WAY OUT OF THE USE OF THIS
% SOFTWARE, EVEN IF ADVISED OF THE POSSIBILITY OF SUCH DAMAGE.

% 此处不用 \specialchap，因为学校要求目录不包括其自己及其之前的内容。
\chapter*{版权声明}
% 综合学校的书面要求及 Word 模版来看，版权声明页不用加页眉、页脚。
\thispagestyle{empty}

任何收存和保管本论文各种版本的单位和个人，
未经本论文作者同意，不得将本论文转借他人，
亦不得随意复制、抄录、拍照或以任何方式传播。
否则，引起有碍作者著作权之问题，将可能承担法律责任。

% 若须排版二维码，请将二维码图片重命名为“barcode”，
% 转为合适的图片格式，并放在当前目录下，然后去掉下面 2 行的注释。
%\vfill\noindent
%\includegraphics[height = 5em]{barcode}

% vim:ts=4:sw=4


\cleardoublepage
\pagestyle{plain}
\setcounter{page}{0}
\pagenumbering{Roman}

\begin{cabstract}
  用一段话概述研究的问题、方法与主要结论。摘要应当独立成篇，
不使用引用和未定义的缩写。
\end{cabstract}

\ifblind\begin{beabstract}\else\begin{eabstract}\fi
  Summarize the problem, the method and the main results in one
paragraph. The abstract stands on its own: no citations and no undefined
abbreviations.
\ifblind\end{beabstract}\else\end{eabstract}\fi

\tableofcontents

\mainmatter

\chapter{引言}
\label{chap:intro}

\section{研究背景}
介绍研究问题的背景与意义，并综述已有工作。引用文献时使用
\cite{knuth1984texbook,lamport1994latex}。

\section{本文的主要工作}
列出本文的主要贡献，并说明全文的结构。

\chapter{研究方法}
\label{chap:method}

描述所用方法。需要引用的公式应当编号，例如
\begin{equation}
  \label{eq:objective}
  \min_{\theta} \; \frac{1}{n} \sum_{i=1}^{n} \ell\bigl(f_\theta(x_i), y_i\bigr) .
\end{equation}

\chapter{实验与分析}
\label{chap:experiments}

表~\ref{tab:results} 给出了求解式~\eqref{eq:objective} 的结果。

\begin{table}[!htb]
  \centering
  \caption{示例数据集上的结果（五次运行的均值 $\pm$ 标准差）}
  \label{tab:results}
  \begin{tabular}{lcc}
    \toprule
    方法     & 准确率（\%）     & 时间（s） \\
    \midrule
    基线方法 & 81.2 $\pm$ 0.4 & 12.4      \\
    本文方法 & 84.7 $\pm$ 0.3 & 10.9      \\
    \bottomrule
  \end{tabular}
\end{table}

\chapter{总结与展望}
总结全文的工作，并指出今后的研究方向。

\printbibliography[heading=bibintoc, title={参考文献}]

\appendix

\chapter{补充材料}
附录用于放置正文以外的推导、数据与程序说明。

\backmatter

\ifblind\else
\chapter{致谢}
感谢导师的指导，以及在研究和论文写作中给予帮助的老师与同学。
\fi

% Copyright (c) 2008-2009 solvethis
% Copyright (c) 2010-2017,2021 Casper Ti. Vector
% Copyright (c) 2021 Kurapica
% All rights reserved.
%
% Redistribution and use in source and binary forms, with or without
% modification, are permitted provided that the following conditions are
% met:
%
% * Redistributions of source code must retain the above copyright notice,
%   this list of conditions and the following disclaimer.
% * Redistributions in binary form must reproduce the above copyright
%   notice, this list of conditions and the following disclaimer in the
%   documentation and/or other materials provided with the distribution.
% * Neither the name of Peking University nor the names of its contributors
%   may be used to endorse or promote products derived from this software
%   without specific prior written permission.
%
% THIS SOFTWARE IS PROVIDED BY THE COPYRIGHT HOLDERS AND CONTRIBUTORS "AS
% IS" AND ANY EXPRESS OR IMPLIED WARRANTIES, INCLUDING, BUT NOT LIMITED TO,
% THE IMPLIED WARRANTIES OF MERCHANTABILITY AND FITNESS FOR A PARTICULAR
% PURPOSE ARE DISCLAIMED. IN NO EVENT SHALL THE COPYRIGHT HOLDER OR
% CONTRIBUTORS BE LIABLE FOR ANY DIRECT, INDIRECT, INCIDENTAL, SPECIAL,
% EXEMPLARY, OR CONSEQUENTIAL DAMAGES (INCLUDING, BUT NOT LIMITED TO,
% PROCUREMENT OF SUBSTITUTE GOODS OR SERVICES; LOSS OF USE, DATA, OR
% PROFITS; OR BUSINESS INTERRUPTION) HOWEVER CAUSED AND ON ANY THEORY OF
% LIABILITY, WHETHER IN CONTRACT, STRICT LIABILITY, OR TORT (INCLUDING
% NEGLIGENCE OR OTHERWISE) ARISING IN ANY WAY OUT OF THE USE OF THIS
% SOFTWARE, EVEN IF ADVISED OF THE POSSIBILITY OF SUCH DAMAGE.

{
	\ctexset{section = {
		format+ = {\centering}, beforeskip = {40bp}, afterskip = {15bp}
	}}
	\specialchap{北京大学学位论文原创性声明和使用授权说明}

	% 学校书面要求本页面不要页码，但在给出的 Word 模版中又有页码。
	% 此处以学校书面要求为准。
	\thispagestyle{empty}
	\mbox{}\vspace*{-3em}
	\section*{原创性声明}

	本人郑重声明：
	所呈交的学位论文，是本人在导师的指导下，独立进行研究工作所取得的成果。
	除文中已经注明引用的内容外，
	本论文不含任何其他个人或集体已经发表或撰写过的作品或成果。
	对本文的研究做出重要贡献的个人和集体，均已在文中以明确方式标明。
	本声明的法律结果由本人承担。
	\vskip 1em
	\rightline{%
		论文作者签名：\hspace{5em}%
		日期：\hspace{2em}年\hspace{2em}月\hspace{2em}日%
	}

	\section*{%
		学位论文使用授权说明\\[-0.33em]
		\textmd{\zihao{5}（必须装订在提交学校图书馆的印刷本）}%
	}

	本人完全了解北京大学关于收集、保存、使用学位论文的规定，即：
	\begin{itemize}
		\item 按照学校要求提交学位论文的印刷本和电子版本；
		\item 学校有权保存学位论文的印刷本和电子版，
			并提供目录检索与阅览服务，在校园网上提供服务；
		\item 学校可以采用影印、缩印、数字化或其它复制手段保存论文；
		\item 因某种特殊原因须要延迟发布学位论文电子版，
			授权学校 $\Box$\nobreakspace{}一年 /
			$\Box$\nobreakspace{}两年 /
			$\Box$\nobreakspace{}三年以后，在校园网上全文发布。
	\end{itemize}
	\centerline{（保密论文在解密后遵守此规定）}
	\vskip 1em
	\rightline{%
		论文作者签名：\hspace{5em}导师签名：\hspace{5em}%
		日期：\hspace{2em}年\hspace{2em}月\hspace{2em}日%
	}

	% 若须排版二维码，请将二维码图片重命名为“barcode”，
	% 转为合适的图片格式，并放在当前目录下，然后去掉下面 2 行的注释。
	%\vfill\noindent
	%\includegraphics[height = 5em]{barcode}
}

% vim:ts=4:sw=4


\end{document}
